\chapter{No complexity in best arm identification}
\label{chap:bai}

Section 5 and Appendix D of \cite{D23}. The answer is the arm with the largest mean, and $\cD$ is
the set of vectors of means with a unique best arm.

\begin{theorem}[{\cite[Theorem 8]{D23}}]
  \label{thm:gaussian_bai_no_complexity}
  \lean{Degenne2023Existence.ofReal_mul_log_le_iInf_maxRatio_baiTask_gaussian}
  \leanok
  \uses{cor:corner_lb, lem:oracleDifficulty_mono, lem:oracleDifficulty_staticClass_le,
    lem:isRegular_baiTask, lem:staticDifficulty_eq_charTime, lem:asymptoticRatio_eq_top}
  For best arm identification with Gaussian arms of variance $1$ and every class $\cC \supseteq
  \Csp$, $\inf_{\cA \in \cC_{\infty}} \sup_{\lambda \in \cD} R_{H_{\cC}, \infty}(\cA, \lambda) \ge
  \frac{3}{80} \log K$.
\end{theorem}

\begin{proof}
  \uses{cor:corner_lb, lem:oracleDifficulty_mono, lem:oracleDifficulty_staticClass_le,
    lem:isRegular_baiTask, lem:staticDifficulty_eq_charTime, lem:asymptoticRatio_eq_top}
  $H_{\cC} \le H_{\Csp} \le \Hsp = T^{\star} \le 2 H_{\Delta}$ (Lemmas~\ref{lem:oracleDifficulty_mono},
  \ref{lem:oracleDifficulty_staticClass_le}, \ref{lem:staticDifficulty_eq_charTime} and the
  bound of \cite{GK16} in LMLExplore), so $R_{H_{\cC}} \ge R_{H_{\Delta}} / 2$. Apply
  Corollary~\ref{cor:corner_lb} at $\mu_1 = 0$, $\mu_k = -k \Delta$ with the arms $j = 2, \dots,
  K$ and $\lambda^{(j)}_j = -\mu_j$: $H_{\Delta}(\lambda^{(j)}) \KL(\mu_j, \lambda^{(j)}_j) \le 8 j$,
  and $\frac18 \sum_{j=2}^K \frac1j \ge \frac{3}{40} \log K$. The paper derives the last
  inequality from $\log((K + 1) / 2) \ge \frac35 \log K$, which fails for $K = 2$ (E6).
\end{proof}

\begin{lemma}[{\cite[Lemma 19]{D23}}]
  \label{lem:exp_fam_bai}
  \lean{Degenne2023Existence.inv_staticDifficulty_baiTask_expFamily_two}
  \leanok
  \uses{def:static_difficulty, lem:exp_fam_bai_min, lem:exp_fam_bai_max}
  For two arms in the same exponential family with log-partition function $\phi$, means $\mu_1 \ne
  \mu_2$ and natural parameters $\xi_1, \xi_2$,
  $\Hsp(\mu)^{-1} = \KL\bigl(\frac{\phi(\xi_1) - \phi(\xi_2)}{\xi_1 - \xi_2}, \mu_1\bigr)$.
\end{lemma}

\begin{proof}
  \uses{lem:exp_fam_bai_min, lem:exp_fam_bai_max}
  The value at the maximizer of Lemma~\ref{lem:exp_fam_bai_max} of the cost of
  Lemma~\ref{lem:exp_fam_bai_min}, $d(\xi_1, y) = \KL(\phi'(y), \mu_1)$ with
  $\phi'(y) = (\phi(\xi_1) - \phi(\xi_2)) / (\xi_1 - \xi_2)$.
\end{proof}

\begin{lemma}[{\cite[Lemma 19]{D23}}]
  \label{lem:exp_fam_bai_min}
  \lean{Degenne2023Existence.iInf_cost_baiTask_expFamily_two}
  \leanok
  \uses{def:static_difficulty}
  In the setting of Lemma~\ref{lem:exp_fam_bai}, for every $\omega \in \Simplex$, the infimum
  over $\lambda \in \Alt(\mu)$ of $\omega_1 \KL(\lambda_1, \mu_1) + \omega_2 \KL(\lambda_2, \mu_2)$
  is its value at $\lambda_1 = \lambda_2 = x(\omega) = \phi'(\omega_1 \xi_1 + \omega_2 \xi_2)$.
\end{lemma}

\begin{proof}
  The alternatives reduce to $\lambda_1 = \lambda_2 = x$ between the means. In natural parameters
  $y$, $\omega_1 d(\xi_1, y) + \omega_2 d(\xi_2, y)$, with $d$ the Bregman divergence of $\phi$, has
  derivative $-(\omega_1 (\xi_1 - y) + \omega_2 (\xi_2 - y)) \phi''(y)$, which vanishes at
  $y = \omega_1 \xi_1 + \omega_2 \xi_2$.
\end{proof}

\begin{lemma}[{\cite[Lemma 19]{D23}}]
  \label{lem:exp_fam_bai_max}
  \lean{Degenne2023Existence.isMaxOn_iInf_cost_baiTask_expFamily_two}
  \leanok
  \uses{lem:exp_fam_bai_min}
  In the setting of Lemma~\ref{lem:exp_fam_bai}, with $x^{\star} = (\phi(\xi_1) - \phi(\xi_2)) /
  (\xi_1 - \xi_2)$, the cost of the alternatives is maximal on the simplex at $\omega^{\star}$ with
  $\omega^{\star}_1 = (\phi'^{-1}(x^{\star}) - \xi_2) / (\xi_1 - \xi_2)$.
\end{lemma}

\begin{proof}
  \uses{lem:exp_fam_bai_min}
  At a maximum, the derivative in $\omega_1$ vanishes, which by the stationarity in $y$ reduces
  to $d(\xi_1, y) = d(\xi_2, y)$, that is $\phi'(y) = x^{\star}$, and $\omega_1 = (y - \xi_2) /
  (\xi_1 - \xi_2)$; the cost is concave in $\omega$.
\end{proof}

\begin{lemma}[{\cite[Lemma 20]{D23}}]
  \label{lem:subgaussian_bai}
  \lean{Degenne2023Existence.ofReal_div_le_inv_staticDifficulty_baiTask_two}
  \leanok
  \uses{lem:exp_fam_bai, lem:rate_ge_sq_of_subgaussian}
  If the members of means $\mu_1$ and $\mu_2$ are $\sigma^2$-sub-Gaussian, then
  $\Hsp(\mu)^{-1} \ge \frac{1}{2 \sigma^2 (\xi_1 - \xi_2)^2} \max\{\KL(\mu_1, \mu_2)^2, \KL(\mu_2,
  \mu_1)^2\}$.
\end{lemma}

\begin{proof}
  \uses{lem:exp_fam_bai, lem:rate_ge_sq_of_subgaussian}
  $x^{\star} = \mu_1 - \KL(\mu_1, \mu_2) / (\xi_1 - \xi_2) = \mu_2 + \KL(\mu_2, \mu_1) / (\xi_1 -
  \xi_2)$, $\Hsp(\mu)^{-1} = \KL(x^{\star}, \mu_1) = \KL(x^{\star}, \mu_2)$ (Lemma~\ref{lem:exp_fam_bai}),
  and $\KL(x, \mu_k) \ge (x - \mu_k)^2 / (2 \sigma^2)$ (Lemma~\ref{lem:rate_ge_sq_of_subgaussian}).
\end{proof}

\begin{lemma}[{\cite[Lemma 20]{D23}}]
  \label{lem:gaussian_bai_two}
  \lean{Degenne2023Existence.inv_staticDifficulty_baiTask_gaussian_two}
  \leanok
  \uses{lem:exp_fam_bai}
  For two Gaussian arms with a common variance $\sigma^2$, the bound of
  Lemma~\ref{lem:subgaussian_bai} is an equality and its two terms are equal; both are
  $(\mu_1 - \mu_2)^2 / (8 \sigma^2)$.
\end{lemma}

\begin{proof}
  \uses{lem:exp_fam_bai}
  $\phi(a) = \frac12 \sigma^2 a^2$, $\KL(x, y) = (x - y)^2 / (2 \sigma^2)$ and
  Lemma~\ref{lem:exp_fam_bai}.
\end{proof}

\begin{lemma}[{\cite[Lemma 9]{D23}}]
  \label{cor:bernoulli_bai_H}
  \lean{Degenne2023Existence.inv_staticDifficulty_baiTask_bernoulli_two}
  \leanok
  \uses{lem:exp_fam_bai}
  For two Bernoulli arms with means $\mu_1 \ne \mu_2$ in $(0, 1)$, with
  $z = \log\frac{1 - \mu_2}{1 - \mu_1} \big/ \log\frac{\mu_1 (1 - \mu_2)}{(1 - \mu_1) \mu_2}$,
  $\Hsp(\mu)^{-1} = \kl(z, \mu_1) = \kl(z, \mu_2)$.
\end{lemma}

\begin{proof}
  \uses{lem:exp_fam_bai}
  Lemma~\ref{lem:exp_fam_bai} with $\phi(a) = \log(1 + e^a)$, $\phi'^{-1}(x) = \log \frac{x}{1 -
  x}$ and the Bernoulli family of LMLExplore.
\end{proof}

\begin{theorem}[{\cite[Theorem 10]{D23}}]
  \label{thm:bernoulli_bai_no_complexity}
  \lean{Degenne2023Existence.one_lt_iInf_maxRatio_baiTask_bernoulli_two}
  \leanok
  \uses{cor:corner_lb, cor:bernoulli_bai_H, lem:oracleDifficulty_mono,
    lem:oracleDifficulty_staticClass_le, lem:isRegular_baiTask, lem:bernoulli_limit,
    lem:asymptoticRatio_eq_top}
  For best arm identification with two Bernoulli arms and every class $\cC \supseteq \Csp$,
  $\inf_{\cA \in \cC_{\infty}} \sup_{\lambda \in \cD} R_{H_{\cC}, \infty}(\cA, \lambda) > 1$.
\end{theorem}

\begin{proof}
  \uses{cor:corner_lb, cor:bernoulli_bai_H, lem:oracleDifficulty_mono,
    lem:oracleDifficulty_staticClass_le, lem:isRegular_baiTask, lem:bernoulli_limit,
    lem:asymptoticRatio_eq_top}
  $H_{\cC} \le \Hsp$ as in Theorem~\ref{thm:gaussian_bai_no_complexity}. For small $x > 0$ (below
  $(\sqrt 3 - 1)/2$, E7), apply Corollary~\ref{cor:corner_lb} at $\mu(x) = (x(1 + x), x)$ with
  $\lambda^{(1)} = (x/2, x)$ and $\lambda^{(2)} = (x(1 + x), 1/2)$; with Lemma~\ref{cor:bernoulli_bai_H}
  the bound tends to about $1.22$ as $x \to 0$ (Lemma~\ref{lem:bernoulli_limit}).
\end{proof}

\begin{lemma}
  \label{lem:bernoulli_limit}
  \uses{cor:bernoulli_bai_H}
  As $x \to 0$,
  \[ \frac{1}{\Hsp((x/2, x)) \kl(x(1 + x), x/2)} + \frac{1}{\Hsp((x(1 + x), 1/2)) \kl(x, 1/2)}
    \to 1 + \frac{1 - \frac{1}{2 \log 2} - \frac{\log(2 \log 2)}{2 \log 2}}{\log 2 - \frac12} >
    1.2. \]
\end{lemma}

\begin{proof}
  \uses{cor:bernoulli_bai_H}
  $\Hsp((x(1 + x), 1/2)) \to 1 / \log 2$ and $\kl(x, 1/2) \to \log 2$; for the first term, both
  $\Hsp((y/2, y))^{-1}$ and $\kl(x, y/2)$ are of order $y$ with $x / y \to 1$, by the formula of
  Lemma~\ref{cor:bernoulli_bai_H}. Numerically the sum exceeds $1$ from $x \approx 10^{-9}$.
\end{proof}

\begin{lemma}
  \label{lem:isRegular_baiTask}
  \uses{asm:answer_nhds}
  Best arm identification on the vectors of means of a family satisfies Assumption 1, and so do
  the positivity and half-space tasks.
\end{lemma}

\begin{proof}
  The answer sets $\{\mu \mid \mu_a > \mu_b \text{ for } b \ne a\}$ are open and regular, and their
  closures cover all vectors.
\end{proof}
